% Part: incompleteness
% Chapter: representability-in-q
% Section: minimization-representable

\documentclass[../../../include/open-logic-section]{subfiles}

\begin{document}

\olfileid[fr]{inc}{req}{min}
\olsection{La minimisation régulière est représentable dans $\Th{Q}$}

Considérons la recherche non bornée. Supposons que $g(x,z)$ soit
totale, régulière et représentable dans~$\Th{Q}$, par la
!!{formula}~$!A_g(x,z,y)$. Définissons
$f(z)=\umin{x}{[g(x,z)=0]}$. Nous cherchons une
!!{formula}~$!A_f(z,y)$ représentant~$f$. La valeur de~$f(z)$
est l'entier~$x$ tel que (a) $g(x,z)=0$ et (b) aucun entier
plus petit ne possède cette propriété : pour tout $w<x$,
$g(w,z)\neq0$. Le choix suivant est donc naturel :
\[
!A_f(z,y) \ident !A_g(y, z, \Obj 0) \land \lforall[w][(w < y \lif \lnot
  !A_g(w, z, \Obj 0))].
\]
Dans le cas général, il faudrait remplacer $z$ par
$z_0,\dots,z_k$.\footnote{Les variables liées de $!A_g$ sont renommées avant les substitutions pour éviter toute capture ; $w$ est choisie distincte des variables d’entrée et de sortie.}

La démonstration exige encore quelques lemmes sur ce que
$\Th{Q}$ est capable de démontrer.

\begin{lem}
\ollabel{lem:succ}
Pour toute !!{constant}~$a$ et tout entier naturel~$n$,
\[
\Th{Q} \Proves \eq[(a' + \num n)][(a + \num n)'].
\]
\end{lem}

\begin{proof}
Raisonnons, comme d'habitude, par récurrence sur~$n$. Pour
$n=0$, il faut démontrer dans~$\Th{Q}$ l'égalité
$\eq[(a'+\Obj 0)][(a+\Obj 0)']$. Or
\begin{align}
  \Th{Q} & \Proves \eq[(a' + \Obj 0)][a'] \quad \text{par l'axiome $Q_4$}
  \ollabel{step1}\\
  \Th{Q} & \Proves  \eq[(a + \Obj 0)][a] \quad \text{par l'axiome $Q_4$}
  \ollabel{step2} \\
  \Th{Q} & \Proves \eq[(a + \Obj 0)'][a'] \quad \text{d'après \olref{step2}}
  \ollabel{step3} \\
  \Th{Q} & \Proves \eq[(a' + \Obj 0)][(a + \Obj 0)'] \quad
  \text{d'après \olref{step1} et \olref{step3}}\notag
\end{align}
Supposons maintenant
$\Th{Q}\Proves\eq[(a'+\num n)][(a+\num n)']$ et
montrons que
$\Th{Q}\Proves\eq[(a'+\num{n}')][(a+\num{n}')']$,
puisque $\num{n+1}=\num n'$. Nous obtenons
\begin{align}
  \Th{Q} & \Proves \eq[(a' + \num n')][(a' + \num n)'] \quad
  \text{par l'axiome $!Q_5$} \ollabel{step5}\\
  \Th{Q} & \Proves \eq[(a' + \num n)'][(a + \num n')'] \quad
  \text{par récurrence, substitution et $!Q_5$} \ollabel{step6}\\
  \Th{Q} & \Proves \eq[(a' + \num n')][(a + \num n')'] \quad
  \text{d'après \olref{step5} et \olref{step6}.} \notag
\end{align}
Pour la deuxième ligne, on applique le successeur à l'égalité
de l'hypothèse de récurrence, puis $!Q_5$ à
$a+\num n'$ : les deux membres obtenus valent
$((a+\num n)')'$.
\end{proof}

Précisons que ce lemme est plus faible que
$\Th{Q}\Proves\lforall[x][\lforall[y][(x'+y)=(x+y)']]$.
Cet !!{sentence} est vrai dans~$\Struct{N}$, mais
$\Th{Q}$ ne le démontre pas.

\begin{lem}
\ollabel{lem:less-zero}
$\Th{Q} \Proves \lforall[x][\lnot x < \Obj 0]$.
\end{lem}

\begin{proof}
Nous raisonnons informellement, en indiquant seulement comment
construire la !!{derivation} formelle.

Pour une !!{constant} arbitraire~$a$, il faut démontrer
$\lnot a<\Obj 0$. D'après la définition de~$<$, il suffit de
démontrer $\lnot\lexists[y][\eq[(y'+a)][\Obj 0]]$ dans~$\Th{Q}$.
Supposons le contraire et prenons un témoin~$b$ tel que
$\eq[(b'+a)][\Obj 0]$. L'axiome~$!Q_3$ donne
$\eq[a][\Obj 0]\lor\lexists[y][\eq[a][y']]$.
Distinguons ces deux cas.

Premier cas : $\eq[a][\Obj 0]$. Alors
$\eq[(b'+\Obj 0)][\Obj 0]$. L'axiome~$!Q_4$ donne
$\eq[(b'+\Obj 0)][b']$, donc $\eq[b'][\Obj 0]$.
Mais $!Q_2$ donne aussi $\eq/[b'][\Obj 0]$ : contradiction.

Second cas : $\eq[a][c']$ pour un certain~$c$. Nous avons
alors $\eq[(b'+c')][\Obj 0]$. D'après $!Q_5$,
$\eq[(b'+c)'][\Obj 0]$, ce qui contredit encore $!Q_2$.
\end{proof}

\begin{lem}
\ollabel{lem:less-nsucc}
Pour tout entier naturel~$n$,
\[
  \Th{Q} \Proves
  \lforall[x][(x < \num {n+1} \lif (\eq[x][\Obj 0] \lor \dots \lor
    \eq[x][\num n]))].
\]
\end{lem}

\begin{proof}
Procédons par récurrence sur~$n$. Pour $n=0$, il faut
montrer $a<\num 1\lif\eq[a][\Obj 0]$, avec~$a$
arbitraire. Supposons d'abord $a<\num 1$. Par l'axiome
définissant~$<$, nous obtenons
$\lexists[y][\eq[(y'+a)][\Obj 0']]$, puisque
$\num 1\ident\Obj 0'$.

Soit $b$ un témoin : $\eq[(b'+a)][\Obj 0']$.
L'axiome~$!Q_3$ donne soit $\eq[a][\Obj 0]$, et il n'y a
rien à faire, soit $\eq[a][c']$ pour un certain~$c$. Dans ce
dernier cas, $\eq[(b'+c')][\Obj 0']$. D'après $!Q_5$,
$\eq[(b'+c)'][\Obj 0']$, puis $!Q_1$ donne
$\eq[(b'+c)][\Obj 0]$. L'axiome~$!Q_8$ entraîne alors
$c<\Obj 0$, en contradiction avec \olref{lem:less-zero}.

Pour l'étape de récurrence, supposons établi le cas~$n$ et
montrons le cas~$n+1$. Si $a<\num{n+2}$, $!Q_3$ fournit
soit $\eq[a][\Obj 0]$, soit $\eq[a][b']$ pour un~$b$.
Dans le premier cas, la disjonction
$\eq[a][\Obj 0]\lor\dots\lor\eq[a][\num{n+1}]$
est immédiate. Dans le second, $b'<\num{n+2}$, autrement dit
$b'<\num{n+1}'$. Par $!Q_8$, il existe un~$c$ tel que
$\eq[(c'+b')][\num{n+1}']$. Par $!Q_5$,
$\eq[(c'+b)'][\num{n+1}']$; par $!Q_1$,
$\eq[(c'+b)][\num{n+1}]$. L'axiome~$!Q_8$ donne
$b<\num{n+1}$. L'hypothèse de récurrence fournit
$\eq[b][\Obj 0]\lor\dots\lor\eq[b][\num n]$.
En appliquant le successeur et $\eq[a][b']$, nous obtenons
$\eq[a][\num 1]\lor\dots\lor\eq[a][\num{n+1}]$.

On peut aussi établir la réciproque. Pour chaque entier fixé $j$ avec
$0\leq j\leq n$, posons $d=n-j$. Le
\olref{lem:succ}, puis le lemme d'addition
\olref[inc][req][bre]{lem:q-proves-add}, donnent
$\Th{Q}\Proves\eq[(\num d'+\num j)][\num{n+1}]$.
L'axiome~$!Q_8$ fournit donc
$\Th{Q}\Proves\num j<\num{n+1}$.
Si $a=\num j$, la substitution des égaux donne
$a<\num{n+1}$. Chaque membre de la disjonction implique
ainsi l'inégalité voulue, ce qui établit l'autre sens de
l'équivalence. Ce raisonnement est une récurrence extérieure
à~$\Th{Q}$ : pour chaque $n$, la disjonction est finie.
\end{proof}

\begin{lem}
\ollabel{lem:trichotomy}
Pour tout entier naturel~$m$,
\[
  \Th{Q} \Proves
  \lforall[y][((y < \num{m} \lor \num{m} < y) \lor \eq[y][\num{m}])].
\]
\end{lem}

\begin{proof}
Raisonnons par récurrence sur~$m$. Pour $m=0$,
l'axiome~$!Q_3$ donne
$\Th{Q}\Proves\lforall[y][(\eq[y][\Obj 0]\lor
\lexists[z][\eq[y][z']])]$.
Soit $a$ arbitraire. Si $\eq[a][\Obj 0]$, la disjonction
voulue suit immédiatement. Sinon, $\eq[a][b']$ pour un~$b$.
La logique de l'égalité donne alors
$\eq[(b'+\Obj 0)][(a+\Obj 0)]$. Par $!Q_4$,
$\eq[(a+\Obj 0)][a]$, d'où
$\eq[(b'+\Obj 0)][a]$ et
$\lexists[z][\eq[(z'+\Obj 0)][a]]$.
L'axiome~$!Q_8$ donne $\Obj 0<a$, donc la disjonction.

Supposons maintenant
\begin{align*}
  \Th{Q} & \Proves \lforall[y][((y < \num{m} \lor \num{m} < y) \lor
    \eq[y][\num{m}])]
  \intertext{et montrons que}
  \Th{Q} & \Proves \lforall[y][((y < \num{m+1} \lor \num{m+1} < y) \lor
    \eq[y][\num{m+1}])]
\end{align*}
Soit $a$ arbitraire. Par $!Q_3$, ou bien
$\eq[a][\Obj 0]$, ou bien $\eq[a][b']$ pour un~$b$.
Dans le premier cas, $!Q_4$ donne
$\eq[\num m'+a][\num{m+1}]$, et $!Q_8$ donne
$a<\num{m+1}$.

Dans le second cas, l'hypothèse de récurrence donne
$(b<\num m\lor\num m<b)\lor\eq[b][\num m]$.
Examinons les trois possibilités.

Si $b<\num m$, alors, d'après $!Q_8$, il existe un~$c$
tel que $\eq[(c'+b)][\num m]$. En appliquant le
successeur, nous avons $\eq[(c'+b)'][\num m']$.
L'axiome~$!Q_5$ donne
$\eq[(c'+b')][(c'+b)']$; ainsi
$\eq[(c'+b')][\num{m+1}]$. En prenant $c$ pour témoin
dans~$!Q_8$, nous obtenons $b'<\num{m+1}$, puis
$a<\num{m+1}$.

Si $\num m<b$, $!Q_8$ fournit un~$c$ tel que
$\eq[(c'+\num m)][b]$. En appliquant le successeur,
$\eq[(c'+\num m)'][b']$. L'axiome~$!Q_5$ donne
$\eq[(c'+\num m')][b']$. Puisque $\eq[a][b']$ et
$\num m'=\num{m+1}$, il en résulte
$\eq[(c'+\num{m+1})][a]$, donc $\num{m+1}<a$
par $!Q_8$.

Enfin, si $\eq[b][\num m]$, le successeur donne
$\eq[b'][\num m']$, donc $\eq[a][\num{m+1}]$.
Chaque membre de la disjonction au rang~$m$ conduit donc au
membre correspondant au rang~$m+1$.
\end{proof}

\begin{prop}
\ollabel{prop:rep-minimization}
Si $g$ est totale et régulière et si $!A_g(x,z,y)$ représente
$g(x,z)$ dans~$\Th{Q}$, alors
\[
!A_f(z,y) \ident !A_g(y, z, \Obj 0) \land \lforall[w][(w < y \lif \lnot
  !A_g(w, z, \Obj 0))]
\]
représente $f(z)=\umin{x}{[g(x,z)=0]}$.
\end{prop}

\begin{proof}
Montrons d'abord que, si $f(n)=m$, alors
$\Th{Q}\Proves !A_f(\num n,\num m)$, c'est-à-dire
\begin{align}
  \Th{Q} & \Proves !A_g(\num{m}, \num{n}, \Obj 0) \land \lforall[w][(w <
    \num{m} \lif \lnot !A_g(w, \num{n}, \Obj 0))]. \notag
  \intertext{Puisque $!A_g$ représente $g$ et que $g(m,n)=0$, nous avons}
\Th{Q} & \Proves !A_g(\num{m}, \num{n}, \Obj 0). \notag
\intertext{Pour chaque $k<m$, la minimalité de $m$ donne $g(k,n)\neq0$, donc}
\Th{Q} & \Proves \lnot !A_g(\num{k}, \num{n}, \Obj 0). \notag
\intertext{Nous en déduisons}
\Th{Q} & \Proves \lforall[w][(w < \num{m} \lif \lnot
  !A_g(w, \num{n}, \Obj 0))]. \ollabel{rep-less}
\end{align}
Pour justifier la troisième ligne, la condition d'unicité de
la représentation de~$g$ donne, à l'entrée $(k,n)$,
$!A_g(\num k,\num n,\Obj 0)\lif\eq[\Obj 0][\num{g(k,n)}]$.
Le \olref[inc][req][bre]{lem:q-proves-neq} fournit la
négation de cette dernière égalité, puisque $g(k,n)\neq0$.
La dernière ligne découle du \olref{lem:less-zero} si $m=0$,
et du \olref{lem:less-nsucc} sinon.

Montrons ensuite que, si $f(n)=m$, alors
$\Th{Q}\Proves\lforall[y][(!A_f(\num n,y)\lif
\eq[y][\num m])]$. Nous esquissons de nouveau la preuve
informellement et laissons au lecteur sa formalisation.

Supposons $!A_f(\num n,b)$. Il en résulte (a)
$!A_g(b,\num n,\Obj 0)$ et (b)
$\lforall[w][(w<b\lif\lnot !A_g(w,\num n,\Obj 0))]$.
Par le \olref{lem:trichotomy},
$(b<\num m\lor\num m<b)\lor\eq[b][\num m]$.
Montrons que chacune des deux inégalités conduit à une
contradiction.

Si $\num m<b$, (b) donne
$\lnot !A_g(\num m,\num n,\Obj 0)$.
Mais $m=f(n)$ implique $g(m,n)=0$ et, comme $!A_g$
représente~$g$, nous avons
$\Th{Q}\Proves !A_g(\num m,\num n,\Obj 0)$ : contradiction.

Si $b<\num m$, la ligne \olref{rep-less} donne
$\lnot !A_g(b,\num n,\Obj 0)$, en contradiction avec~(a).
Il reste donc $\eq[b][\num m]$, comme voulu.
\end{proof}

\begin{explain}
La régularité de~$g$ assure qu'à chaque entrée correspond
un premier zéro. Pour vérifier l'instance positive de la
représentation, le \olref{lem:less-nsucc} ramène les valeurs
strictement inférieures à un numéral fixé à une disjonction
finie. Pour vérifier l'unicité, le \olref{lem:trichotomy}
permet de comparer une valeur candidate à ce numéral,
même lorsque la candidate est une variable du langage.
\end{explain}

\end{document}
